\section{TRUY VẤN BẰNG ĐẠI SỐ QUAN HỆ}
\label{sec:chuong4}

Để biểu thức gọn hơn, nhóm viết tắt tên các quan hệ như sau: TDP = to\_dan\_pho, DC = dia\_chi, HO = ho\_cu\_tru, CD = cu\_dan, QT = qua\_trinh\_cu\_tru, TK = tai\_khoan và BD = bien\_dong\_cu\_tru. Nhiều bảng cùng có cột ghi\_chu, nên nếu dùng phép kết nối tự nhiên (natural join) thì cột này cũng bị đem ra so sánh và cho kết quả sai. Vì vậy, ở các câu có kết nối, nhóm hoặc chiếu bớt cột trước khi kết nối, hoặc dùng phép kết nối theo điều kiện ($\theta$-join) ghi rõ cột cần so khớp.

\phantomsection
\label{subsec:qd}
Các câu hỏi về người đang cư trú xét hiệu lực theo quy tắc B05 tại ngày hiện tại của phiên PostgreSQL. Đại số quan hệ cổ điển không có NULL, nên nhóm chuẩn bị đầu vào bằng SQL trước khi áp dụng các phép toán. Quan hệ $Q_{\mathrm{ht}}$ gồm các đợt còn hiệu lực và chỉ giữ các cột bắt buộc:
\begin{equation*}
\begin{aligned}
Q_{\mathrm{ht}}(&\text{ma\_qua\_trinh},\ \text{ma\_cu\_dan},\ \text{ma\_ho},\\
    &\text{loai\_cu\_tru},\ \text{ngay\_bat\_dau})
\end{aligned}
\end{equation*}
Năm cột trên đều NOT NULL. Quan hệ $Q_{\mathrm{bd}}$ có cùng lược đồ nhưng gồm tất cả đợt đã bắt đầu. Bước chuẩn bị $Q_{\mathrm{bd}}$ lọc \texttt{ngay\_bat\_dau <= CURRENT\_DATE}; $Q_{\mathrm{ht}}$ còn yêu cầu ngày kết thúc để trống hoặc không trước \texttt{CURRENT\_DATE}. Cả hai dùng SELECT DISTINCT năm cột trên, không đưa cột ngày kết thúc vào đầu vào đại số. Cách xử lý này loại đợt chưa bắt đầu và giữ NULL ở ngoài các phép toán cổ điển.

\textbf{Câu 1.} Cho biết họ tên và ngày sinh của các cư dân nữ, phục vụ việc lập danh sách hoạt động của chi hội phụ nữ.
\begin{equation*}
\pi_{\text{ho\_ten},\ \text{ngay\_sinh}}\left(\sigma_{\text{gioi\_tinh} = \text{'Nu'}}(CD)\right)
\end{equation*}
Phép chọn $\sigma$ lọc ra các bộ có giới tính là Nu, sau đó phép chiếu $\pi$ chỉ giữ lại hai cột cần in. Nhóm đặt phép chọn ở trong và phép chiếu ở ngoài, vì nếu chiếu trước thì cột gioi\_tinh đã bị bỏ đi và không còn để lọc nữa.

\textbf{Câu 2.} Liệt kê mã và địa chỉ chuẩn của tất cả địa chỉ thuộc ``Tổ 01''.
\begin{equation*}
\begin{aligned}
R &\leftarrow \sigma_{\text{ten\_to} = \text{'Tổ 01'}}(TDP)\\
KQ &\leftarrow \pi_{\text{ma\_dia\_chi},\ \text{dia\_chi\_chuan}}\left(R \bowtie_{R.\text{ma\_to} = DC.\text{ma\_to}} DC\right)
\end{aligned}
\end{equation*}
Nhóm chọn tổ trước rồi mới kết nối để giảm số bộ ở đầu vào của phép kết nối theo biến đổi đại số. Đây là cách tổ chức biểu thức, không phải kết luận đo hiệu năng thực thi.

\textbf{Câu 3.} Cho biết mã và họ tên các cư dân đang thường trú tại tổ có mã 1. Biểu thức khá dài nên nhóm viết thành từng bước với các quan hệ trung gian.
\begin{equation*}
\begin{aligned}
R_1 &\leftarrow \sigma_{\text{loai\_cu\_tru} = \text{'THUONG\_TRU'}}(Q_{\mathrm{ht}})\\
R_2 &\leftarrow \pi_{\text{ma\_ho}}\left(HO \bowtie_{HO.\text{ma\_dia\_chi} = DC.\text{ma\_dia\_chi}} \sigma_{\text{ma\_to} = 1}(DC)\right)\\
R_3 &\leftarrow R_1 \bowtie_{R_1.\text{ma\_ho} = R_2.\text{ma\_ho}} R_2\\
KQ &\leftarrow \pi_{CD.\text{ma\_cu\_dan},\ \text{ho\_ten}}\left(R_3 \bowtie_{R_3.\text{ma\_cu\_dan} = CD.\text{ma\_cu\_dan}} CD\right)
\end{aligned}
\end{equation*}
$R_1$ là các đợt thường trú còn hiệu lực hiện tại, $R_2$ là các hộ nằm trong tổ 1; $R_3$ ghép hai kết quả này lại, và kết nối $R_3$ với CD là lấy được họ tên.

\textbf{Câu 4.} Tìm các cư dân đã từng cư trú trong phường nhưng hiện tại không còn đợt nào hiệu lực, để cán bộ rà soát lại hồ sơ.
\begin{equation*}
\begin{aligned}
R &\leftarrow \pi_{\text{ma\_cu\_dan}}\left(Q_{\mathrm{bd}}\right) - \pi_{\text{ma\_cu\_dan}}(Q_{\mathrm{ht}})\\
KQ &\leftarrow \pi_{\text{ma\_cu\_dan},\ \text{ho\_ten}}(R \bowtie CD)
\end{aligned}
\end{equation*}
Phép hiệu lấy những người có đợt đã bắt đầu trong $Q_{\mathrm{bd}}$, trừ đi những người còn đợt hiệu lực trong $Q_{\mathrm{ht}}$. Cả hai vế đều chiếu về mã cư dân, nên tương thích với phép hiệu. Ở bước sau, $R$ chỉ có cột ma\_cu\_dan nên dùng kết nối tự nhiên với CD là an toàn.

\textbf{Câu 5.} Cho biết mã các cư dân đã từng tạm trú và cũng đã từng thường trú trong phường.
\begin{equation*}
\begin{aligned}
A &\leftarrow \pi_{\text{ma\_cu\_dan}}\left(\sigma_{\text{loai\_cu\_tru} = \text{'TAM\_TRU'}}(QT)\right)\\
B &\leftarrow \pi_{\text{ma\_cu\_dan}}\left(\sigma_{\text{loai\_cu\_tru} = \text{'THUONG\_TRU'}}(QT)\right)\\
KQ &\leftarrow A \cap B
\end{aligned}
\end{equation*}
Không thể viết câu này bằng một phép chọn với điều kiện AND, vì mỗi bộ của QT chỉ có một loại cư trú; hai điều kiện phải đúng ở hai bộ khác nhau nên cần dùng phép giao.

\textbf{Câu 6.} Liệt kê mã tất cả các biến động có liên quan đến cư dân mã 92, dù cư dân này nằm ở đợt cũ hay đợt mới của biến động.
\begin{equation*}
\begin{aligned}
Q &\leftarrow \pi_{\text{ma\_qua\_trinh}}\left(\sigma_{\text{ma\_cu\_dan} = 92}(QT)\right)\\
KQ &\leftarrow \pi_{\text{ma\_bien\_dong}}\left(BD \bowtie_{\text{ma\_qua\_trinh\_cu} = \text{ma\_qua\_trinh}} Q\right)\\
&\phantom{\leftarrow}\ \cup\ \pi_{\text{ma\_bien\_dong}}\left(BD \bowtie_{\text{ma\_qua\_trinh\_moi} = \text{ma\_qua\_trinh}} Q\right)
\end{aligned}
\end{equation*}
Vế trái của phép hợp lấy các biến động mà cư dân 92 rời đi, vế phải lấy các biến động mà cư dân 92 đến; một lần chuyển nội bộ sẽ xuất hiện ở cả hai vế nhưng phép hợp tự loại bỏ bộ trùng.

\textbf{Câu 7.} Với mỗi biến động chuyển nội bộ, cho biết họ tên cư dân, mã hộ cũ và mã hộ mới.
\begin{equation*}
\begin{aligned}
QC &\leftarrow \rho_{QC(\text{mq\_cu},\ \text{ho\_cu})}\left(\pi_{\text{ma\_qua\_trinh},\ \text{ma\_ho}}(QT)\right)\\
P &\leftarrow \pi_{\text{ma\_qua\_trinh},\ \text{ma\_ho},\ \text{ma\_cu\_dan}}(QT)\\
QM &\leftarrow \rho_{QM(\text{mq\_moi},\ \text{ho\_moi},\ \text{ma\_cu\_dan})}(P)\\
R_0 &\leftarrow \sigma_{\text{loai\_bien\_dong} = \text{'CHUYEN\_NOI\_BO'}}(BD)\\
R_1 &\leftarrow R_0 \bowtie_{\text{ma\_qua\_trinh\_cu} = \text{mq\_cu}} QC\\
R &\leftarrow R_1 \bowtie_{\text{ma\_qua\_trinh\_moi} = \text{mq\_moi}} QM\\
KQ &\leftarrow \pi_{\text{ma\_bien\_dong},\ \text{ho\_ten},\ \text{ho\_cu},\ \text{ho\_moi}}\left(R \bowtie \pi_{\text{ma\_cu\_dan},\ \text{ho\_ten}}(CD)\right)
\end{aligned}
\end{equation*}
Vì một biến động tham chiếu tới QT hai lần với hai vai trò khác nhau, nhóm dùng phép đổi tên $\rho$ để tạo ra hai bản sao QC (đợt cũ) và QM (đợt mới) có tên cột khác nhau, nhờ đó kết nối được cả hai lần mà không bị trùng tên thuộc tính.

\textbf{Câu 8.} Cho biết tên hiển thị của những cán bộ đã ghi nhận biến động chuyển đi xảy ra trong tháng 10/2026.
\begin{equation*}
\begin{aligned}
R &\leftarrow \sigma_{\substack{\text{loai\_bien\_dong} = \text{'CHUYEN\_DI'}\\ \wedge\ \text{ngay\_bien\_dong} \ge \text{'2026-10-01'}\\ \wedge\ \text{ngay\_bien\_dong} \le \text{'2026-10-31'}}}(BD)\\
KQ &\leftarrow \pi_{\text{ten\_hien\_thi}}\left(R \bowtie_{R.\text{ma\_tai\_khoan} = TK.\text{ma\_tai\_khoan}} TK\right)
\end{aligned}
\end{equation*}
Phép chọn trên BD gộp ba điều kiện bằng phép hội $\wedge$ để lấy đúng loại biến động trong khoảng thời gian cần xét, rồi mới kết nối với TK để lấy tên cán bộ. Vì phép chiếu tự loại các bộ trùng nhau, một cán bộ ghi nhiều lần chuyển đi trong tháng cũng chỉ xuất hiện một lần trong kết quả.

\textbf{Câu 9.} Tìm những tài khoản đã ghi nhận đủ các loại biến động xuất hiện trong dữ liệu.
\begin{equation*}
\pi_{\text{ma\_tai\_khoan},\ \text{loai\_bien\_dong}}(BD) \div \pi_{\text{loai\_bien\_dong}}(BD)
\end{equation*}
Quan hệ bị chia $\pi_{\text{ma\_tai\_khoan},\ \text{loai\_bien\_dong}}(BD)$ gồm các cặp (tài khoản, loại biến động) đã từng được ghi; quan hệ chia $\pi_{\text{loai\_bien\_dong}}(BD)$ là tập các loại biến động có trong hệ thống. Kết quả có lược đồ là phần còn lại $\{\text{ma\_tai\_khoan}\}$, gồm những tài khoản mà khi ghép với \emph{mọi} bộ của quan hệ chia đều có mặt trong quan hệ bị chia. Đây là dạng câu hỏi có chữ ``tất cả'' mà các phép chọn, chiếu, kết nối thông thường không diễn đạt trực tiếp được. Quan hệ chia ở đây lấy từ dữ liệu đã ghi; nếu muốn hiểu ``tất cả loại'' theo miền giá trị của CHECK thì thay bằng quan hệ hằng gồm ba bộ CHUYEN\_DEN, CHUYEN\_DI, CHUYEN\_NOI\_BO. Trên bộ dữ liệu mẫu cả ba loại đều đã xuất hiện nên hai cách hiểu cho cùng kết quả. Nhóm kiểm chứng bằng hai câu SQL độc lập, một câu dùng hai lớp NOT EXISTS và một câu đếm số loại phân biệt của mỗi tài khoản: cả hai đều trả về bốn tài khoản 1, 2, 3, 4; tài khoản 5 chưa ghi biến động nào nên không có bộ nào trong quan hệ bị chia và không thể xuất hiện trong kết quả.

\textbf{Câu 10.} Đếm số nhân khẩu đang cư trú hiện tại của từng hộ.
\begin{equation*}
{}_{\text{ma\_ho}}\,\gamma_{\,\text{COUNT(ma\_cu\_dan)} \rightarrow \text{so\_nhan\_khau}}\left(\pi_{\text{ma\_ho},\ \text{ma\_cu\_dan}}(Q_{\mathrm{ht}})\right)
\end{equation*}
Phép gom nhóm $\gamma$ (một số giáo trình ký hiệu là $\Im$) không thuộc các phép toán của đại số quan hệ cổ điển mà thuộc \textbf{đại số quan hệ mở rộng}, vốn bổ sung phép gom nhóm và hàm tổng hợp như COUNT \cite{elmasri,silberschatz}. Phép này chia các cặp (hộ, cư dân) theo ma\_ho rồi đếm số cư dân trong mỗi nhóm. Nhóm chiếu về cặp (ma\_ho, ma\_cu\_dan) trước khi gom nhóm để một người có cả đợt thường trú và đợt tạm trú trong cùng hộ chỉ được đếm một lần, đúng với quy tắc B17 là số người phải tính theo mã cư dân phân biệt. Trên bộ dữ liệu mẫu hiện tại chưa có hộ nào rơi vào trường hợp này, nhưng phép chiếu vẫn cần để biểu thức đúng với mọi dữ liệu hợp lệ. Kết quả có 56 hộ đang có người hiện tại, trong đó 40 hộ có từ 3 nhân khẩu trở lên, khớp với câu 3 ở Chương \ref{sec:chuong5}.

\subsection*{Đối chiếu biểu thức đại số quan hệ với SQL}

Mỗi biểu thức được dịch sang SQL theo ngữ nghĩa tập hợp: phép chiếu dùng SELECT DISTINCT; hợp, giao, hiệu lần lượt dùng UNION, INTERSECT, EXCEPT; phép chia dùng hai lớp NOT EXISTS; phép gom nhóm dùng GROUP BY. Các câu tương ứng ở Chương \ref{sec:chuong5} được so sánh tập kết quả theo cả hai chiều. Bảng \ref{tab:doichieura} trình bày số liệu minh họa trên bộ dữ liệu mẫu tại thời điểm thực nghiệm. Các kết quả về cư trú hiện tại có thể thay đổi theo ngày hệ thống; việc kiểm chứng đối chiếu truy vấn dùng cùng ngày hệ thống, không gắn kết quả với một ngày cố định.

\begin{table}[H]
\centering
\fontsize{11pt}{14pt}\selectfont
\renewcommand{\arraystretch}{1.15}
\begin{tabularx}{\textwidth}{|P{1.3cm}|P{2.4cm}|Y|}
\hline
\textbf{Câu} & \textbf{Số bộ kết quả} & \textbf{Đối chiếu} \\
\hline
1 & 98 & Bằng số cư dân nữ (không có hai người nữ trùng cả họ tên và ngày sinh) \\
\hline
2 & 8 & Bằng số địa chỉ thuộc Tổ 01 \\
\hline
3 & 23 & Trùng khớp từng bộ với câu 1 Chương \ref{sec:chuong5} \\
\hline
4 & 17 & Trùng khớp từng bộ với câu 4 Chương \ref{sec:chuong5} \\
\hline
5 & 31 & Bằng số dòng của câu 9 Chương \ref{sec:chuong5}, vốn dùng cùng phép INTERSECT \\
\hline
6 & 1 (biến động 32) & Lần chuyển nội bộ ngày 02/7/2026 từ hộ 56 sang hộ 20; trùng với cách viết một phép chọn dùng điều kiện ``hoặc'' \\
\hline
7 & 27 & Trùng khớp từng bộ (4 cột) với câu 7 Chương \ref{sec:chuong5} \\
\hline
8 & 2 & Cán bộ Cư trú 01 và Cán bộ Cư trú 02, từ 2 biến động chuyển đi trong tháng 10/2026 \\
\hline
9 & 4 (tài khoản 1--4) & Cách viết NOT EXISTS và cách đếm cho cùng kết quả \\
\hline
10 & 56 hộ & 40 hộ có từ 3 nhân khẩu, khớp câu 3 Chương \ref{sec:chuong5} \\
\hline
\end{tabularx}
\caption{Kết quả đối chiếu 10 câu đại số quan hệ trên bộ dữ liệu mẫu}
\label{tab:doichieura}
\end{table}
